In this section we present an effective algorithm to compute the girth of a conjugacy class in a free group.  This gives an algorithm to compute the girth of a closed curve in a handlebody $H$.  To compute the girth of a curve, simply pass to the level of the fundamental group $\pi_{1}(H)$.  Using the results of this section, we also show that free groups of rank bigger than one do not contain elements of girth one.  We prove this at the end.

In the previous section we introduced the notion of a $\dd$-interface.  Let $H$ be a handlebody.  Let $\dd$ be a disc system for $H$.  Let $\gamma$ be a closed curve in $H$, and let $\dint{a}$ be a $\dd$-interface.  We define the intersection number $\gamma\wedge\dint{a}$ as follows.  Lift $\dint{a}$ to an interface $\alpha$ for $\hh$, as in Section~\ref{subordination}.  This lift is well-defined up to deck transformations.  Consider all lifts of $\gamma$ to $\hh$.  Define the intersection number $\gamma\wedge\dint{a}$ to be the number of lifts of $\gamma$ that cross the interface $\alpha$.  By equivariance, this quantity does not depend on the particular lift of $\dint{a}$.  The following lemma says how we can compute the intersection number.

\begin{lemma}
\label{intersection handlebody}
Let $H$ be a handlebody.  Let $\dd$ be a disc system for $H$.  Let $\gamma$ be a closed curve in $H$.  Let $\dint{a}$ be a $\dd$-interface.  Suppose $\gamma$ minimizes intersection with $\dd$ in its homotopy class.  The disc system $\dd$ divides $\gamma$ into a collection of arcs connecting elements of $\dpm$.  The intersection number $\gamma\wedge\dint{a}$ equals the number of arcs of $\gamma$ connecting a black element of $\dpm$ to a white element.
\end{lemma}

\begin{proof}
Let $g$ be the genus of $H$.  Form the directed graph $\Gamma_{\dd}$ consisting of a single vertex and $g$ directed edges in correspondence with $\dd$.  There exists a homotopy $\Phi:H\rightarrow\Gamma_{\dd}$.  Choose $\Phi$ so that the preimage under $\Phi$ of the midpoints of the edges in $\Gamma_{\dd}$ equals $\dd$.

Consider the universal covers of $H$ and $\Gamma_{\dd}$.  Denote them $\hh$ and $\Gamma$.  The homotopy $\Phi$ lifts to a homotopy from $\hh$ to $\Gamma$.  The graph $\Gamma$ is the tree associated to $\dd$.  We can also view $\Gamma$ as the Cayley graph of $\pi_{1}(H)$ with respect to the generating system corresponding to $\dd$.  Lift $\gamma$ to a path in $\Gamma$.  By Lemma~\ref{geodesic}, the curve $\gamma$ lifts to a path without backtrackings.

Choose a vertex $v$ of $\Gamma$ and lift the $\dd$-interface $\dint{a}$ to an interface $\alpha$ for $\Gamma$.  The $2g$ edges that touch $v$ are in correspondence with $\dpm$.  Each of these $2g$ edges inherits a color black, white, or gray from $\dint{a}$.  The intersection number $\gamma\wedge\dint{a}$ equals the number of lifts of $\gamma$ to $\Gamma$ that cross $\alpha$.  Any geodesic path in $\Gamma$ that crosses $\alpha$ must pass through a black edge touching $v$, the vertex $v$, and a white edge touching $v$.  

Consider all lifts of $\gamma$ to $\Gamma$ that pass through $v$.  These lifts are in correspondence with the arcs $\gamma-D$ in $H$.  The lifts crossing $\alpha$ correspond to the arcs connecting a black element of $\dpm$ to a white element, so we are done.
\end{proof}

Let $F$ be a free group.  Let $\dd$ be a generating system for $F$.  Let $w$ be an element of $F$, and let $\dint{a}$ be a $\dd$-interface.  Then $\dint{a}$ assigns a color of black, white, or gray to each of the semigroup generators in $\dpm$.  In parallel with the above, we define the intersection number $w\wedge\dint{a}$.  Form a handlebody $H$ with disc system $\dd$.  Choose a basepoint for $H$.  The group $\pi_{1}(H)$ is naturally in correspondence with $F$.  Choose a closed loop $\gamma$ in $H$ that gives rise to $w$.  View $\dint{a}$ as a $\dd$-interface for the handlebody $H$.  Define $w\wedge\dint{a}$ as $\gamma\wedge\dint{a}$.  This definition does not depend on the choice of $H$, the choice of closed curve, or the choice of basepoint.  The following lemma is simply Lemma~\ref{intersection handlebody} phrased in terms of free groups.

\begin{lemma}
\label{intersection free group}
Let $F$ be a free group.  Let $\dd$ be a generating system for $F$.  Let $w$ be a word in the semigroup generators $\dpm$ of $F$.  Let $\dint{a}$ be a $\dd$-interface.  Suppose that $w$ is cyclically reduced.  The intersection number $w\wedge\dint{a}$ equals the number of adjacent pairs $ab$ in $w$ such that $a^{-1}$ is black and $b$ is white, or $a^{-1}$ is white and $b$ is black.
\end{lemma}

\begin{proof}
Realize $F$ as the fundamental group of a handlebody $H$ with disc system $\dd$, as above.  Consider any closed curve $\gamma$ in $H$ that is transverse to the discs in $\dd$.  By looking at the sequence of intersections with elements of $\dd$, the curve $\gamma$ gives rise to a word in the elements of $\dpm$.  In other words, $\gamma$ gives rise to a word in the generators $\dd$ and their inverses.

Let $\Gamma$ denote the tree associated to $\dd$.  Choose a closed curve $\gamma$ in $H$ that gives rise to the word $w$.  Since $w$ is cyclically reduced, the curve $\gamma$ has no backtrackings when lifted to $\Gamma$.  By Lemma~\ref{geodesic}, the curve $\gamma$ minimizes intersection with $\dd$ in its homotopy class.

The quantity $w\wedge\dint{a}$ equals $\gamma\wedge\dint{a}$ by definition.  By Lemma~\ref{intersection handlebody}, the quantity $\gamma\wedge\dint{a}$ equals the number of arcs of $\gamma-D$ connecting a black element of $\dpm$ to a white element.   An arc of $\gamma-D$ corresponds to an adjacent pair of letters in the word $w$.  The arcs of $\gamma-D$ that connect a black element of $\dpm$ to a white element correspond to the adjacent pairs of letters $ab$ such that $a^{-1}$ is black, resp. white, and $b$ is white, resp. black.
\end{proof}

Let $w$ be an element of the free group $F$.  Let $\dd$ be a disc system for $F$.  Define the \emph{complexity} of $\dd$ with respect to $w$ to be $w\wedge\dd$.  By Lemma~\ref{cyclic}, the complexity of $\dd$ equals the length of $w$ written as a cyclically reduced word in the generators corresponding to $\dd$.

Let $\dint{c}$ be a $\dd$-half-disc $\di{\db}{\dw}$.  Let $\dg$ the corresponding gray disc.  We say $\dint{c}$ is a \emph{shortcut} $\dd$-half-disc for $\dd$ with respect to $w$ if $w\hi\di{\db}{\dw}$ is strictly less than both $\gamma\hi\di{\db}{\dg}$ and $\gamma\hi\di{\dw}{\dg}$.  Note that $\gamma\hi\di{\db}{\dg}$, for instance, is the $\dd$-interface obtained by taking $\dd_{B}$ as the set of black discs and $\dd_{G}$ as the set of white discs.

By the following lemma, if $\dd$ has a shortcut $\dd$-half-disc, we can create a generating system of smaller complexity.  Thus we say $\dd$ is \emph{locally minimal} with respect to $w$ if $\dd$ has no shortcut $\dd$-half-discs.  Locally minimal disc systems are useful for calculating girth by Lemma~\ref{locally-minimal}.

\begin{lemma}
\label{curve shortcut}
Let $F$ be a free group.  Let $\dd$ be a generating system for $F$.  Let $w$ be an element of $F$.  If there is a shortcut $\dd$-half-disc with respect to $w$, then there is a generating system of smaller complexity with respect to $w$.
\end{lemma}

\begin{proof}
The set $\dpm$ generates $F$ as a semigroup.  Suppose $\di{\db}{\dw}$ is a shortcut $\dd$-half-disc with respect to $w$, with gray element $\dg$.   Then $w\hi\di{\db}{\dw}<\min(w\hi\di{\db}{\dg},w\hi\di{\dw}{\dg})$.  The inverse of the element corresponding to $\dg$ in $\dpm$ is either black or white.  Assume the inverse is black, say, so that it lies in $\db$.  We now work geometrically.

Realize the free group $F$ and generating system $\dd$ as the fundamental group of a handlebody $H$ with disc system labeled by $\dd$.  Choose a closed curve $\gamma$ that gives rise to $w$.  Then $\gamma\hi\di{\db}{\dw}<\min(\gamma\hi\di{\db}{\dg},\gamma\hi\di{\dw}{\dg})$.

Let $D_{G}$ denote the geometric disc in $\dd$ that corresponds to $\dg$.  Realize the $\dd$-disc $\di{\db}{\dg\cup\dw}$ as a new geometric disc $D_{N}$ in $H-\dd$.  Replace $D_{G}$ with $D_{N}$ to create a new disc system $\dd^\prime$.  The collection $\dd^\prime$ is a disc system because, by construction, $D_{N}$ is non-separating in $H$:  the collection $\db$ contains the element of $\dpm$ paired with $\dg$ in $H$.

The complexity of $\dd^{\prime}$ is the sum of $\gamma\hi D$ over all discs $D$ in $\dd$.  The discs in $\dd^{\prime}$ are the same as those in $\dd$, except that $D_{N}$ replaced $D_{G}$.  These two discs satisfy $\gamma\hi D_{N}=\gamma\hi\di{\db}{\dg}+\gamma\hi\di{\db}{\dw}<\gamma\hi\di{\db}{\dg}+\gamma\hi\di{\dw}{\dg}=\gamma\hi D_{G}$.  The middle inequality follows from the shortcut inequality $\gamma\hi\di{\db}{\dw}<\gamma\hi\di{\dw}{\dg}$.  Since $\gamma\hi D_{N}<\gamma\hi D_{G}$, the disc system $\dd^\prime$ has smaller complexity than $\dd$.
\end{proof}

The above proof is constructive.  It shows how to construct a disc system of smaller complexity if a shortcut $\dd$-half-disc exists.  The construction is geometric.  For algorithmic purposes we restate the construction algebraically below.

Let $F$ be a free group with generating system $\dd$.  The elements of $\dpm$ generate $F$ as a semigroup.  Let $w$ be a word in the generators $\dpm$.  Let $\di{\db}{\dw}$ be a shortcut $\dd$-half-disc with respect to $w$, with gray element $\dg$.  Cyclically reduce $w$.  By Lemma~\ref{cyclic}, the complexity of $\dd$ is the length of $w$.

We will create a new generating system $\dd^{\prime}$ by constructing an automorphism $\Phi$ of $F$.  Automorphisms preserve girth.  By Lemma~\ref{curve shortcut}, the word $\Phi(w)$, when cyclically reduced, will have length shorter than $w$ with respect to the generating system $\dd$.  By taking $\dd^{\prime}$ to be $\Phi^{-1}(\dd)$, we will obtain a generating system of smaller complexity with respect to $w$.

Denote as $c$ the group element in $\dpm$ corresponding to $\dg$.  Switch the black and white subsets of $\di{\db}{\dw}$, if necessary, so that $c^{-1}$ is black.  This mirrors the proof of Lemma~\ref{curve shortcut}.  If we translate the proof of Lemma~\ref{curve shortcut} into algebraic terms, we get the automorphism defined in Table~\ref{automorphism}.

\begin{table}
\begin{center}
\begin{tabular}{c|c|c|}
\multicolumn{3}{r}{black \ \ white} \\ \cline{2-3}
black & $c^{-1}ac$ & \ $c^{-1}a$ \ \\ \cline{2-3}
white & $ac$ & $a$ \\ \cline{2-3 }
\end{tabular}
\end{center}
\caption{This table gives instructions on where $\Phi$ should map each generator $a$ of the generating system $\dd$ for $F$.  The row corresponds to whether $a$ is black or white as an element of $\dpm$.  The column corresponds to whether $a^{-1}$ is black or white as an element of $\dpm$.}
\label{automorphism}
\end{table}



\begin{lemma}
\label{locally-minimal}
Let $F$ be a free group.  Let $\dd$ be a generating system for $F$.  Let $w$ be an element of $F$.  If $\dd$ is locally minimal with respect to $w$, then the girth of $w$ equals the minimum of $w\hi\dint{a}$ over all $\dd$-discs $\dint{a}$.
\end{lemma}

\begin{proof}
Form a handlebody $H$ with disc system labeled by $\dd$.  Choose a basepoint for $H$.  The fundamental group $\pi_{1}(H)$ is in natural correspondence with $F$.  Let $\gamma$ be a closed loop in $H$ that gives rise to $w$.  Any $\dd$-disc for $F$ transfers to a $\dd$-disc for $H$.  If $\dint{a}$ is a $\dd$-disc for $F$, then $w\hi\dint{a}$ equals $\gamma\hi\dint{a}$, with the second expression taken in $H$.  By definition, the girth of $w$ is the minimum of $\gamma\hi D$ over all essential discs $D$ in $H$.  Let $n$ be the girth of $\gamma$.  Let $D$ be a girth-realizing disc for $\gamma$, so that $\gamma\hi D=n$.  Apply Lemma~\ref{subordination-lemma} to the disc $D$ in $H$ with disc system $\dd$.  There are two cases.

\textbf{Case 1:  There is a $\dd$-disc subordinate to $D$.}
Let $\dint{a}$ be a $\dd$-disc subordinate to $D$.  Since $\dint{a}$ 
is subordinate to $D$, we have $\gamma\hi\dint{a}\leq\gamma\hi D$.
Thus $\gamma\hi\dint{a}\leq n$.  Proceed with $\dint{a}$ to the conclusion 
below.

\textbf{Case 2:  There are two $\dd$-half-discs independently 
subordinate to $D$.}  Since there are two $\dd$-half-discs independently subordinate to $D$, one of the $\dd$-half-discs, call it $\dint{c}$, satisfies 
$\gamma\hi\dint{c}\leq\lfloor\frac{n}{2}\rfloor$.  Let $\di{\db}{\dw}$ denote $\dint{c}$, with gray disc $\dg$.  Then $\gamma\hi\di{\db}{\dw}\leq\lfloor\frac{n}{2}\rfloor$.  Since $\dd$ is locally minimal, the $\dd$-half-disc $\dint{c}$ is not shortcut.  So $\gamma\hi\di{\db}{\dw}\geq\gamma\hi\di{\db}{\dg}$, say.  Let $\dint{a}$ be the $\dd$-disc $\di{\db}{\dg\cup\dw}$.  Then $\gamma\hi\dint{a}$ equals $\gamma\hi\di{\db}{\dg}+\gamma\hi\di{\db}{\dw}\leq n$.  Proceed to the conclusion.

\textbf{Conclusion of Cases.}
In both cases we have a $\dd$-disc $\dint{a}$ satisfying 
$\gamma\hi\dint{a}\leq n$.  Consider a $\dd$-disc $\dint{b}$ that minimizes $w\hi\dint{b}$.  Since $\dint{a}$ exists, we know $w\hi\dint{b}\leq n$.  Realize $\dint{b}$ by a geometric disc $D$ in $H$.  Then $\gamma\hi D=w\hi\dint{b}$.  Since $n$ is the girth of $\gamma$, we have $\gamma\hi D\geq n$.  Therefore 
$w\hi\dint{b}=n$.
\end{proof}

Now we present our algorithm to compute the girth of an element of a free group.  The algorithm is exponential in the rank of the free group and linear in the length of the initial word in the generators and their inverses.

\textbf{The Algorithm.}  Let $F$ be a free group of rank $g$ with generating system $\dd$.  Let $w$ be an element of $F$ represented as a word in the elements of $\dpm$.  We compute the girth of the conjugacy class of $w$.

\textbf{Step 1:  Find a shortcut $\dd$-half-disc with respect to $w$, if possible.}  There are finitely many $\dd$-half-discs.  Indeed, there are $2n(2^{2n-2}-1)$ $\dd$-half-discs, up to switching black and white.  For each $\dd$-half-disc, use Lemma~\ref{intersection free group} to check whether it is a shortcut $\dd$-half-disc with respect to $w$.  If there is a shortcut $\dd$-half-disc, apply Lemma~\ref{curve shortcut}.  Replace $\dd$ with a generating system of smaller complexity, and repeat \textbf{Step~1}.  If there is no shortcut $\dd$-half-disc, then $\dd$ is locally minimal.  Go to \textbf{Step~2}.

\textbf{Step 2:  Find a $\dd$-disc $\dint{a}$ that minimizes $w\hi\dint{a}$.}
There are finitely many $\dd$-discs.  Indeed, there are $2^{2n-1}-1$ $\dd$-discs, up to switching black and white.  For each $\dd$-disc $\dint{a}$, use Lemma~\ref{intersection free group} to calculate $w\hi\dint{a}$.  Find one that minimizes $w\hi\dint{a}$.  There may be several.  By Lemma~\ref{locally-minimal}, the quantity $w\hi\dint{a}$ equals the girth of $w$.

\textbf{End Algorithm.}

Using the results of this section, we prove here that free groups of rank two or more do not contain elements of girth one.  This gives an alternative proof of Proposition~\ref{algebraic girth one}.

\begin{lemma}
\label{group girth one}
Let $F$ be a free group.  If $F$ has rank bigger than one, then $F$ does not contain an element of girth one.
\end{lemma}

\begin{proof}
Let $F$ have rank two or more.  Let $w$ be an element of $F$ with $\girth(w)\leq 1$.  Realize $F$ as the fundamental group of a handlebody $H$.  Represent $w$ as a closed curve $\gamma$ in $H$.  Then the girth of $w$ equals the girth of $\gamma$.

Let $D$ be a disc such that $\gamma\wedge D\leq 1$.  If $\gamma\wedge D=0$, then $\gamma$ has girth zero, and we are done.  So assume $\gamma\wedge D=1$.  Homotope $\gamma$ in $H$ to minimize intersection with $D$.  By Lemma~\ref{wedge}, the curve $\gamma$ intersects $D$ exactly once.  Since $\gamma$ intersects $D$ an odd number of times, the disc $D$ is non-separating.  Complete $D$ to a disc system $\dd$.  To finish the proof, it suffices to construct a disc system of smaller complexity or show that $\gamma$ has girth zero.

Homotope $\gamma$ to minimize intersection with $\dd$.  By Corollary~\ref{subcollection}, the curve $\gamma$ continues to intersect $D$ once.  Let $D_{B}$ be one of the two oriented discs in $\dpm$ that correspond to the disc $D$.  Let $\kappa$ be the arc component of $\gamma-\dd$ that leaves $D_{B}$.  Let $D_{G}$ be the disc in $\dpm$ at the other end of $\kappa$.  Let $\di{\db}{\dw}$ denote the $\dd$-half-disc defined by taking $D_{G}$ as the gray disc $\dg$ and $D_{B}$ as the collection of black discs $\db$.  Let $\dw$ denote the white discs, which make up the rest of $\dpm$.

By construction, we have $w\wedge\di{\db}{\dw}=0$ and $w\wedge\di{\dg}{\db}=1$.  If $w\wedge\di{\dg}{\dw}=0$, then $\di{\dw}{\dg\cup\db}$ defines a $\dd$-disc $\dint{b}$ with $w\wedge\dint{b}=0$.  Realizing $\dint{b}$ as a geometric disc, we see that $w$ would have girth zero.  So assume that $w\wedge\di{\dg}{\dw}>0$.  In this case, the $\dd$-half-disc $\di{\db}{\dw}$ is a shortcut $\dd$-half-disc.  By Lemma~\ref{curve shortcut}, we can reduce the complexity of $\dd$, so we are done.
\end{proof}